\chapter{Background and Related Work}
\label{ch:background}

This chapter surveys the bodies of work that inform AI Compass's design: personalized recommendation, semantic retrieval over dense embeddings, result diversification, retrieval-augmented generation, ranked-retrieval evaluation, and speech recognition for long-form video. Each section closes by relating the literature to a specific, concrete choice made in the system, rather than treating the two as separate concerns. The final section synthesizes what the system adopts from this literature and what it deliberately does not, and why.

\section{Personalized Recommendation Systems}
\label{sec:background-recsys}

Personalized recommendation is conventionally divided into content-based filtering and collaborative filtering, with hybrid approaches combining the two \cite{ricci2015recommender}. Content-based filtering represents items by their own features (text, topic, metadata) and represents a user by an aggregate of the features of items that user has previously engaged with, then recommends items whose feature representation is close to the user's profile \cite{pazzani2007content}. Collaborative filtering instead represents a user by their pattern of engagement relative to other users, and recommends items that similar users engaged with, without necessarily modeling item content at all \cite{su2009collaborative}. Collaborative approaches (including matrix factorization and, more recently, learned two-tower models) are attractive because they can surface non-obvious item-to-item relationships purely from co-engagement, but they depend structurally on having enough users and enough overlapping engagement history for meaningful co-engagement patterns to emerge. A new platform, or one with a small user base, faces the classic \emph{cold-start} problem: with few users, there is not yet a dense enough interaction matrix for collaborative signal to be distinguishable from noise.

AI Compass adopts a content-based, feature-driven approach rather than a collaborative one, and this is a deliberate, documented scope decision rather than a missing feature. The project's own roadmap records the reasoning explicitly: collaborative filtering, matrix factorization, and two-tower architectures are postponed because, at a user count on the order of tens rather than thousands, cross-user co-engagement signal is cold-start by construction, and the content-based recommender is judged the right approach at the platform's present scale, with collaborative approaches to be revisited only once meaningful monthly active usage accumulates. Every signal the ranking service consumes (interest match, content quality, freshness, source affinity, and novelty relative to what a user has already been shown, formalized fully in the Personalization and Recommendation chapter (Chapter~6)) is computed from an individual user's own profile and behavior, never from other users' co-engagement patterns. Table~\ref{tab:cb-vs-cf} summarizes the trade-off that motivates this choice.

\begin{table}[htbp]
  \centering
  \caption{Content-based versus collaborative filtering, and the basis for AI Compass's choice.}
  \label{tab:cb-vs-cf}
  \begin{tabular}{|p{3.2cm}|p{6cm}|p{6cm}|}
    \hline
    \textbf{Property} & \textbf{Content-based filtering} & \textbf{Collaborative filtering} \\
    \hline
    Signal source & Item features + own history & Cross-user co-engagement \\
    \hline
    Cold-start (new user) & Degrades gracefully (feature match) & Fails without history \\
    \hline
    Cold-start (small platform) & Usable from day one & Requires many users \\
    \hline
    Explainability & Direct (feature-level) & Indirect (latent factors) \\
    \hline
    AI Compass's choice & Adopted & Deferred until scale justifies it \\
    \hline
  \end{tabular}
\end{table}

\section{Semantic Retrieval and Embeddings}
\label{sec:background-embeddings}

Modern semantic retrieval rests on the Transformer architecture \cite{vaswani2017attention}, whose self-attention mechanism allows a model to weigh every token in a sequence against every other token, producing contextual representations that outperformed prior recurrent architectures on sequence-modeling tasks. Sentence-BERT adapted Transformer encoders such as BERT to produce fixed-length sentence and passage embeddings that are directly comparable by cosine similarity, avoiding the need to run the full model as a cross-encoder over every candidate pair at query time \cite{reimers2019sentencebert}. This made embedding-based semantic search practical at scale: passages are embedded once, offline, and a query embedding is compared against a large precomputed index rather than being jointly encoded with every candidate.

AI Compass embeds content with \texttt{all-MiniLM-L6-v2}, a compact sentence-transformers model in the MiniLM family. MiniLM itself is a knowledge-distillation method for compressing large pretrained Transformers into smaller student models that retain most of the teacher's task performance at a fraction of the inference cost \cite{wang2020minilm}. It is important to state this lineage precisely rather than overclaim it: \texttt{all-MiniLM-L6-v2} is a sentence-transformers checkpoint distributed and documented through the sentence-transformers project and its Hugging Face model card \cite{sbert_docs}, and its lineage to the original MiniLM distillation method runs through that model-card documentation rather than through a direct, one-to-one correspondence with the original MiniLM paper's own trained model. The distinction matters for a thesis that is committed to citation integrity: the paper describes and validates the distillation method that the model family is built on, not the specific 384-dimensional checkpoint AI Compass loads.

Once content is embedded, finding the nearest neighbors of a query vector in a large collection by exhaustive comparison does not scale. Hierarchical Navigable Small World (HNSW) graphs address this by building a multi-layer proximity graph over the indexed vectors, so that an approximate nearest-neighbor search can be performed by greedily traversing the graph from a coarse layer down to a fine one, trading a small, bounded amount of recall for large gains in query speed \cite{malkov2020hnsw}. AI Compass indexes its embedding columns with pgvector, an extension that adds a vector column type and both exact and HNSW approximate-nearest-neighbor index types directly to PostgreSQL \cite{pgvector}. Running the vector index inside the same relational database that already stores every other piece of content metadata avoids introducing and operationally maintaining a second, specialized vector-database service, at the cost of depending on pgvector's specific index implementation and tuning parameters rather than a purpose-built vector engine.

\section{Diversification and Result Ranking}
\label{sec:background-mmr}

A ranking function that only ever selects the single highest-scoring next item tends to produce a result list dominated by near-duplicates of whichever topic scores best, which is precisely the filter-bubble failure mode a personalization system should avoid. Maximal Marginal Relevance (MMR) addresses this by reformulating selection as a greedy trade-off: at each step, the next item chosen is the one that maximizes a linear combination of its own relevance score and a penalty for its similarity to items already selected, controlled by a single interpolation parameter \cite{carbonell1998mmr}. The method was originally proposed for reordering retrieved documents and for extractive summarization, but it generalizes directly to any setting where a ranked list should balance individual relevance against list-level diversity. AI Compass's ranking service applies exactly this selection rule over its scored candidate pool; the specific interpolation weight the project uses, and the additional bounded exploration mechanism layered on top of it, are project-specific configuration choices covered in full in the Personalization and Recommendation chapter (Chapter~6), not properties of the MMR method itself.

\section{Retrieval-Augmented Generation}
\label{sec:background-rag}

Retrieval-Augmented Generation (RAG) combines a parametric generative language model with a non-parametric retrieval step over an external corpus: rather than relying solely on facts encoded in the model's weights at training time, the model is conditioned at generation time on passages retrieved for the specific query, which both grounds its output in verifiable source material and allows the corpus to be updated without retraining the model \cite{lewis2020rag}. This is the foundational framing behind any system, including AI Compass's conversational assistant, that answers questions by first retrieving relevant passages and then generating an answer conditioned on them.

Dense embedding-based retrieval, as described in Section~\ref{sec:background-embeddings}, is not the only retrieval strategy available, and it is not infallible as an infrastructure dependency: if the embedding service that produces query and passage vectors is unavailable, a system built purely on dense retrieval has no fallback. The classical alternative is sparse, term-based retrieval, exemplified by the BM25 ranking function, which scores documents against a query by a weighted combination of term frequency and inverse document frequency without requiring any learned representation at all \cite{robertson2009bm25}. AI Compass's search module falls back to exactly this kind of keyword-based retrieval when the embedding service is unavailable, trading semantic recall for a dependency-free retrieval path that degrades gracefully rather than failing outright.

\section{Evaluation Metrics for Ranked Retrieval}
\label{sec:background-eval}

Evaluating a ranked list of recommendations or search results requires metrics that account for both the relevance of returned items and their position in the list, since a relevant item near the top of a list is more valuable to a user than the same item buried near the bottom. Normalized Discounted Cumulative Gain (NDCG) formalizes this by discounting each item's relevance contribution by a function of its rank position, then normalizing against the score of an ideal ordering so that scores are comparable across queries with different numbers of relevant items \cite{jarvelin2002ndcg}. Mean Average Precision (MAP) instead averages precision computed at each point a relevant item appears in the ranked list, summarizing how consistently relevant items are front-loaded rather than scattered throughout the list; both metrics, along with the underlying definitions of precision and recall they build on, are given a standard treatment in the information-retrieval literature \cite{manning2008ir}. AI Compass's offline evaluation harness computes NDCG@k and MAP against held-out behavioral events, providing the quantitative basis for the ranking-quality discussion in the Evaluation chapter.

\section{Speech Recognition for Long-Form Video}
\label{sec:background-asr}

Long-form video is one of the source types AI Compass ingests, and video whose host platform does not supply a caption track requires automatic transcription before its content can be summarized, embedded, or searched at all. Whisper is a Transformer-based automatic speech recognition model trained on a very large and weakly supervised corpus of audio paired with transcripts drawn from the web, which gives it substantially more robust performance across accents, background noise, and technical vocabulary than models trained on smaller, more curated speech corpora \cite{radford2023whisper}. AI Compass's speech-to-text stage does not run the reference Whisper implementation directly; it runs \texttt{faster-whisper}, a reimplementation of Whisper inference on top of the CTranslate2 inference engine that reproduces the same model architecture and published weights while optimizing the inference runtime itself, which is what makes CPU-only transcription at usable throughput feasible without dedicated GPU hardware.

\section{Positioning of AI Compass Relative to Prior Work}
\label{sec:background-positioning}

AI Compass adopts, largely unmodified, four ideas from the literature surveyed above: Maximal Marginal Relevance for diversifying its ranked feed \cite{carbonell1998mmr}, HNSW-backed approximate nearest-neighbor search for its embedding-based candidate retrieval and passage search \cite{malkov2020hnsw}, the retrieval-then-generate framing of retrieval-augmented generation for its conversational assistant \cite{lewis2020rag}, and NDCG/MAP as its offline ranking-quality metrics \cite{jarvelin2002ndcg,manning2008ir}. Each of these is used as published, with project-specific parameter choices (an interpolation weight, a candidate pool size, an evaluation window) rather than a modification to the underlying method.

Equally deliberately, AI Compass does not adopt three approaches that a broader reading of the literature might suggest. It does not implement collaborative filtering (Section~\ref{sec:background-recsys}), because the project's own roadmap documents that the platform's current user count puts any cross-user co-engagement signal squarely in cold-start territory, and revisiting collaborative approaches is explicitly deferred until usage grows. It does not use a learned or neural ranking model in its serving path: ranking is a fixed, documented, deterministic weighted formula rather than a model trained on logged features, consistent with the project's architecture principle that the ranking hot path stays deterministic and free of both machine-learned scoring models and large-language-model calls, so that every production ranking decision is reproducible and explainable without re-running inference. It does not fine-tune either the embedding model or the generative language models it calls; both are used as distributed, off-the-shelf checkpoints \cite{sbert_docs,groq_docs}. These are not omissions discovered after the fact but positions taken up front for reproducibility and data-scale reasons: a system that ingests content once and personalizes at read time (the project's governing "ingest once, personalize later" principle) has no natural per-item retraining hook, and a serving path that must stay fast, cheap, and explainable at every request has no room for either a generative model or a continuously retrained ranker sitting on its hot path. Where machine learning is used at all, it is used offline and once per item, in the enrichment step described in the Data Acquisition and Content Intelligence chapter (Chapter~5), not in the request-serving path itself.
